% Aula 08 --- a ROC só depende da ordem dos escores; o corte de custo depende da escala.
Carla é epidemiologista e avalia um modelo que estima, para cada paciente, a probabilidade
de uma doença ($Y=1$). Numa amostra de teste com $4$ doentes e $6$ saudáveis, independente
dos dados usados para ajustar o modelo, ela anota o escore $s_i$ do seu modelo e o escore
$r_i$ do modelo de um colega:
\begin{center}\small
\renewcommand{\arraystretch}{1.2}
\begin{tabular}{@{}lcccccccccc@{}}
\toprule
paciente & 1 & 2 & 3 & 4 & 5 & 6 & 7 & 8 & 9 & 10 \\
\midrule
$y_i$ & 1 & 1 & 1 & 1 & 0 & 0 & 0 & 0 & 0 & 0 \\
$s_i$ & $0{,}90$ & $0{,}70$ & $0{,}45$ & $0{,}25$ & $0{,}60$ & $0{,}35$ & $0{,}15$ & $0{,}12$ & $0{,}08$ & $0{,}05$ \\
$r_i$ & $0{,}988$ & $0{,}845$ & $0{,}401$ & $0{,}100$ & $0{,}692$ & $0{,}225$ & $0{,}030$ & $0{,}018$ & $0{,}008$ & $0{,}003$ \\
\bottomrule
\end{tabular}
\end{center}
Ela percebe que, em todos os pacientes, $r_i=h(s_i)$, com $h(s)=s^2/\big(s^2+(1-s)^2\big)$.
Cada modelo declara positivo quem tem escore maior ou igual a um corte $p_0$, e a curva ROC
de um escore é o conjunto dos pontos (taxa de falsos positivos, taxa de verdadeiros
positivos) que os cortes produzem nesta amostra, ligados por segmentos.
\begin{itens}
  \item \pts{1,2} Liste os pontos da curva ROC do escore $s$, desenhe a curva e calcule a
        área sob ela, a AUC. Em que ponto da curva fica o corte $p_0=0{,}5$?
  \item \pts{1,0} Mostre que $h$ é estritamente crescente em $(0,1)$ e que, para
        \emph{toda} função $h$ estritamente crescente, o escore $h(s)$ tem a mesma curva
        ROC que $s$ nesta amostra. Conclua que os dois modelos têm a mesma AUC.
  \item \pts{1,0} No hospital, deixar passar um doente custa quatro vezes mais que um
        alarme falso: $\ell_0=4$ e $\ell_1=1$. O corte que minimiza o custo esperado, quando
        o escore é a probabilidade verdadeira $\Prob(Y=1\mid\X=\x)$, é
        $\ell_1/(\ell_0+\ell_1)$. Aplique esse corte aos escores $s$ e aos escores $r$,
        identifique os pacientes declarados positivos em cada caso e calcule o custo total
        $\ell_1\,\mathrm{FP}+\ell_0\,\mathrm{FN}$ de cada decisão na amostra. Compare com o
        custo dos outros cortes possíveis para $s$.
  \item \pts{0,8} O que a curva ROC e a AUC dizem sobre os dois modelos, e o que elas não
        dizem? O que é preciso saber sobre os escores para usar o corte
        $\ell_1/(\ell_0+\ell_1)$?
\end{itens}

\begin{solucao}
\textbf{(a)} Os escores são distintos. Em ordem decrescente, as classes aparecem na
sequência $1,1,0,1,0,1,0,0,0,0$ (pacientes 1, 2, 5, 3, 6, 4, 7, 8, 9, 10). Baixar o corte
de um escore ao seguinte acrescenta um paciente aos declarados positivos: se ele é doente,
a taxa de verdadeiros positivos sobe $1/P=1/4$; se é saudável, a de falsos positivos sobe
$1/N=1/6$. Começando sem ninguém declarado positivo, os pontos são
\[
  (0,0),\ \big(0,\tfrac14\big),\ \big(0,\tfrac12\big),\ \big(\tfrac16,\tfrac12\big),\
  \big(\tfrac16,\tfrac34\big),\ \big(\tfrac13,\tfrac34\big),\ \big(\tfrac13,1\big),\
  \big(\tfrac12,1\big),\ \big(\tfrac23,1\big),\ \big(\tfrac56,1\big),\ (1,1),
\]
e a curva é uma escada: sobe nos doentes e anda para a direita nos saudáveis. A área sob
ela é a soma das áreas dos retângulos sob os trechos horizontais, cada um de largura $1/6$:
\[
  \mathrm{AUC}=\tfrac16\Big(\tfrac12+\tfrac34+1+1+1+1\Big)=\tfrac16\cdot\tfrac{21}{4}=\tfrac78=0{,}875 .
\]
O corte $0{,}5$ declara positivos os pacientes 1, 2 e 5: dois doentes e uma saudável, o ponto
$\big(\tfrac16,\tfrac12\big)$.

\smallskip
\textbf{(b)} Em $(0,1)$, $h(s)=1/\big(1+((1-s)/s)^2\big)$. A função $s\mapsto(1-s)/s=1/s-1$ é
estritamente decrescente e positiva, o seu quadrado também é estritamente decrescente, e
$u\mapsto1/(1+u)$ é estritamente decrescente em $u>0$; a composta é estritamente crescente.

Seja agora $h$ uma função estritamente crescente qualquer. Para todo corte $p_0$ do escore
$s$, o conjunto dos declarados positivos é $\{i: s_i\ge p_0\}$, e, como $h$ é estritamente
crescente, $s_i\ge p_0\iff h(s_i)\ge h(p_0)$. Esse conjunto é, portanto, exatamente o que o
corte $h(p_0)$ produz para o escore $h(s)$. Reciprocamente, cada corte $q$ do escore $h(s)$
declara positivos os pacientes com $h(s_i)\ge q$, que formam o conjunto
$\{i: s_i\ge s_{(q)}\}$, com $s_{(q)}$ o menor $s_i$ tal que $h(s_i)\ge q$, ou nenhum paciente,
se não há tal $s_i$. Os dois escores produzem, assim, a mesma família de conjuntos de
declarados positivos, e cada conjunto determina as duas taxas. As curvas ROC têm os mesmos
pontos e, ligadas por segmentos, são a mesma curva, com a mesma área: os dois modelos têm
AUC $0{,}875$.

\smallskip
\textbf{(c)} O corte é $\ell_1/(\ell_0+\ell_1)=1/5=0{,}2$. Nos escores $s$, ficam acima dele
os pacientes 1, 2, 3, 4, 5 e 6: os quatro doentes e duas saudáveis, com $\mathrm{FN}=0$,
$\mathrm{FP}=2$ e custo $1\cdot2+4\cdot0=2$. Nos escores $r$, ficam acima de $0{,}2$ os
pacientes 1, 2, 3, 5 e 6, e o paciente 4, doente, com $r_4=0{,}100$, fica de fora: $\mathrm{FN}=1$,
$\mathrm{FP}=2$ e custo $1\cdot2+4\cdot1=6$. Os outros cortes de $s$ custam, do mais alto ao
mais baixo, $16$ (ninguém positivo), $12$, $8$, $9$, $5$, $6$, $2$, $3$, $4$, $5$ e $6$
(todos positivos): o corte $0{,}2$ aplicado a $s$ escolhe, nesta amostra, o ponto mais
barato da curva, $\big(\tfrac13,1\big)$, e aplicado a $r$ escolhe o ponto
$\big(\tfrac13,\tfrac34\big)$, que custa três vezes mais.

A razão está na escala. Como $\log\frac{h(s)}{1-h(s)}=2\log\frac{s}{1-s}$, o escore $r$
dobra o \emph{log-odds} de $s$, e $r>0{,}2$ equivale a $s>1/3$: o mesmo número $0{,}2$, aplicado
a escalas diferentes, seleciona pontos diferentes da mesma curva. Em $10$ pacientes, o custo
observado é uma estimativa ruidosa do custo esperado; o que o exercício mostra é o mecanismo,
não o tamanho da diferença.

\smallskip
\textbf{(d)} A curva ROC e a AUC dizem que trocas entre as duas taxas cada escore permite,
corte a corte, e isso só depende da ordem em que ele põe os pacientes; nisso os dois modelos
são idênticos. Não dizem nada sobre a escala dos escores, porque qualquer transformação
estritamente crescente as deixa iguais: uma curva ROC não distingue uma probabilidade bem
estimada de outra que exagere a certeza. O corte $\ell_1/(\ell_0+\ell_1)$, por sua vez, é ótimo para a
probabilidade verdadeira $\Prob(Y=1\mid\X=\x)$; aplicado a um escore, ele só tem esse sentido
se o escore estimar bem essa probabilidade. Se Carla tem razões para confiar nas suas
probabilidades, o modelo do colega, que dobra o \emph{log-odds}, exagera a certeza, como faz o
Bayes ingênuo com covariáveis redundantes, e o corte $0{,}2$ aplicado a $r$ deixa passar
doentes. Para escolher o corte, é preciso saber se os escores estão na escala de
probabilidades; para escolher entre modelos com a mesma ROC, a curva não basta.
\end{solucao}
